\textbf{Phần 1. Trắc nghiệm lựa chọn}
\setcounterpageref{ex}{0}

\begin{ex} %[3D2V13]
Tìm thừa số $x$, biết: $x \times 6 = 18 + 24$.
\choice
{$6$}
{$8$}
{\True $7$}
{$9$}
\loigiai{
Tính vế phải: $x \times 6 = 42 \Rightarrow x = 42 : 6 = 7$.
\textbf{Đáp án của câu hỏi là C}
}
\end{ex}

\begin{ex} %[3D2V13]
Tìm số bị chia $x$, biết: $x : 7 = 35 : 5$.
\choice
{$42$}
{\True $49$}
{$56$}
{$35$}
\loigiai{
Tính vế phải: $x : 7 = 7 \Rightarrow x = 7 \times 7 = 49$.
\textbf{Đáp án của câu hỏi là B}
}
\end{ex}

\begin{ex} %[3D2V13]
Tìm số chia $x$, biết: $54 : x = 24 : 4$.
\choice
{$8$}
{$7$}
{$6$}
{\True $9$}
\loigiai{
Tính vế phải: $54 : x = 6 \Rightarrow x = 54 : 6 = 9$.
\textbf{Đáp án của câu hỏi là D}
}
\end{ex}

\begin{ex} %[3D2V13]
Trong một phép nhân, nếu một thừa số gấp lên $3$ lần và thừa số kia giữ nguyên thì tích mới thay đổi như thế nào?
\choice
{\True Gấp lên $3$ lần}
{Giảm đi $3$ lần}
{Tăng thêm $3$ đơn vị}
{Không thay đổi}
\loigiai{
Khi nhân một thừa số với 3 thì tích cũng gấp lên 3 lần.
\textbf{Đáp án của câu hỏi là A}
}
\end{ex}

\begin{ex} %[3D2V13]
Tìm $x$ biết: $x \times 8 + 16 = 80$.
\choice
{$7$}
{$9$}
{\True $8$}
{$6$}
\loigiai{
Xem $x \times 8$ là số hạng chưa biết: $x \times 8 = 80 - 16 = 64 \Rightarrow x = 64 : 8 = 8$.
\textbf{Đáp án của câu hỏi là C}
}
\end{ex}

\begin{ex} %[3D2V13]
Tìm $x$ biết: $x : 9 - 3 = 5$.
\choice
{$63$}
{\True $72$}
{$81$}
{$54$}
\loigiai{
Xem $x : 9$ là số bị trừ: $x : 9 = 5 + 3 = 8 \Rightarrow x = 8 \times 9 = 72$.
\textbf{Đáp án của câu hỏi là B}
}
\end{ex}

\begin{ex} %[3D2V13]
Trong một phép chia hết có thương bằng $6$. Nếu số bị chia gấp lên $2$ lần và giữ nguyên số chia thì thương mới là
\choice
{$8$}
{$3$}
{$18$}
{\True $12$}
\loigiai{
Khi số bị chia gấp lên 2 lần thì thương cũng gấp lên 2 lần: $6 \times 2 = 12$.
\textbf{Đáp án của câu hỏi là D}
}
\end{ex}

\begin{ex} %[3D2V13]
Một số nhân với $7$ rồi cộng với $15$ thì được $50$. Số đó là
\choice
{\True $5$}
{$6$}
{$4$}
{$7$}
\loigiai{
Gọi số đó là $x$: $x \times 7 + 15 = 50 \Rightarrow x \times 7 = 35 \Rightarrow x = 5$.
\textbf{Đáp án của câu hỏi là A}
}
\end{ex}

\begin{ex} %[3D2V13]
Tìm $x$ biết: $72 : x + 12 = 20$.
\choice
{$8$}
{$7$}
{\True $9$}
{$6$}
\loigiai{
Ta có: $72 : x = 20 - 12 = 8 \Rightarrow x = 72 : 8 = 9$.
\textbf{Đáp án của câu hỏi là C}
}
\end{ex}

\begin{ex} %[3D2V13]
Tích của hai số bằng $48$. Nếu giữ nguyên thừa số thứ nhất và giảm thừa số thứ hai đi $2$ lần thì tích mới là
\choice
{$16$}
{\True $24$}
{$32$}
{$96$}
\loigiai{
Khi giảm một thừa số đi 2 lần thì tích cũng giảm đi 2 lần: $48 : 2 = 24$.
\textbf{Đáp án của câu hỏi là B}
}
\end{ex}

\begin{ex} %[3D2V13]
Tìm $x$ biết: $(x + 5) \times 7 = 63$.
\choice
{$3$}
{$5$}
{$6$}
{\True $4$}
\loigiai{
Ta có: $x + 5 = 63 : 7 = 9 \Rightarrow x = 9 - 5 = 4$.
\textbf{Đáp án của câu hỏi là D}
}
\end{ex}

\begin{ex} %[3D2V13]
Một thùng chứa có $64\text{ lít}$ nước. Người ta dùng các can loại $8\text{ lít}$ để múc hết nước trong thùng. Hỏi cần ít nhất bao nhiêu chiếc can như vậy?
\choice
{\True $8$ chiếc can}
{$7$ chiếc can}
{$9$ chiếc can}
{$6$ chiếc can}
\loigiai{
Số chiếc can cần dùng là: $64 : 8 = 8$ chiếc can.
\textbf{Đáp án của câu hỏi là A}
}
\end{ex}

\textbf{Phần 2. Trắc nghiệm đúng sai}
\setcounterpageref{ex}{0}

\begin{ex} %[3D2V13]
Xét tính đúng sai của các quy tắc tìm thành phần phép tính:
\choiceTF
{\True Muốn tìm thừa số chưa biết, ta lấy tích chia cho thừa số đã biết.}
{Muốn tìm số chia, ta lấy thương nhân với số bị chia.}
{\True Muốn tìm số bị chia trong phép chia hết, ta lấy thương nhân với số chia.}
{\True Trong phép chia có dư, số bị chia bằng thương nhân số chia cộng số dư.}
\loigiai{
\begin{itemize}
    \item Ý a) Đúng: Quy tắc cơ bản của phép nhân.
    \item Ý b) Sai: Muốn tìm số chia, ta lấy số bị chia chia cho thương.
    \item Ý c) Đúng: Quy tắc tìm số bị chia.
    \item Ý d) Đúng: Công thức phép chia có dư.
\end{itemize}
\begin{itemize}
    \item a) \textbf{Đúng}.
    \item b) \textbf{Sai}.
    \item c) \textbf{Đúng}.
    \item d) \textbf{Đúng}.
\end{itemize}
}
\end{ex}

\begin{ex} %[3D2V13]
Cho bài toán tìm $x$: $x \times 8 - 14 = 50$. Xét tính đúng sai:
\choiceTF
{\True Giá trị của biểu thức $x \times 8$ bằng $64$.}
{Giá trị của $x$ tìm được là $9$.}
{\True Số $x$ tìm được là một số chẵn chia hết cho 4.}
{\True Nếu thay số $50$ thành $58$ thì giá trị của $x$ sẽ bằng $9$.}
\loigiai{
\begin{itemize}
    \item Ý a) Đúng: $50 + 14 = 64$.
    \item Ý b) Sai: $x = 64 : 8 = 8$ (không phải 9).
    \item Ý c) Đúng: $8$ là số chẵn chia hết cho 4.
    \item Ý d) Đúng: $x \times 8 = 58 + 14 = 72 \Rightarrow x = 9$.
\end{itemize}
\begin{itemize}
    \item a) \textbf{Đúng}.
    \item b) \textbf{Sai}.
    \item c) \textbf{Đúng}.
    \item d) \textbf{Đúng}.
\end{itemize}
}
\end{ex}

\textbf{Phần 3. Trả lời ngắn}
\setcounterpageref{ex}{0}

\begin{ex} %[3D2V13]
Tìm thừa số $x$, biết: $x \times 7 = 84 - 28$.
\shortans{8}
\loigiai{
\textbf{Đáp số:} $8$\\
$x \times 7 = 56 \Rightarrow x = 56 : 7 = 8$.
}
\end{ex}

\begin{ex} %[3D2V13]
Tìm số bị chia $x$, biết: $x : 6 = 15 - 8$.
\shortans{42}
\loigiai{
\textbf{Đáp số:} $42$\\
$x : 6 = 7 \Rightarrow x = 7 \times 6 = 42$.
}
\end{ex}

\begin{ex} %[3D2V13]
Tìm số chia $x$, biết: $72 : x = 27 : 3$.
\shortans{8}
\loigiai{
\textbf{Đáp số:} $8$\\
$72 : x = 9 \Rightarrow x = 72 : 9 = 8$.
}
\end{ex}

\begin{ex} %[3D2V13]
Một số nhân với 8 rồi trừ đi 15 thì được 49. Tìm số đó.
\shortans{8}
\loigiai{
\textbf{Đáp số:} $8$\\
$x \times 8 - 15 = 49 \Rightarrow x \times 8 = 64 \Rightarrow x = 8$.
}
\end{ex}

\begin{ex} %[3D2V13]
Tìm $x$ biết: $(x - 6) \times 9 = 54$.
\shortans{12}
\loigiai{
\textbf{Đáp số:} $12$\\
$x - 6 = 54 : 9 = 6 \Rightarrow x = 6 + 6 = 12$.
}
\end{ex}

\begin{ex} %[3D2V13]
Trong một phép chia có thương là 8, số chia là 7 và số dư là 4. Tìm số bị chia.
\shortans{60}
\loigiai{
\textbf{Đáp số:} $60$\\
$8 \times 7 + 4 = 56 + 4 = 60$.
}
\end{ex}
